% !TEX root = ../monograph.tex
\section{Bounded real approximations for separated words}
\label{ch13:bounded-approximations}

The later phase argument needs real approximants bounded uniformly in
the cutoff. Compactness of the link space at each cutoff is insufficient:
the physical normalization of a local field grows as the mesh shrinks.
We instead stop at one fixed coarse level and pull back a bounded coarse
function through the probability block kernels.

We record the precise marked estimate used in this construction. It is
the local-trial comparison for reflected tilings, with every term involving at least two insertions
retained, not only its one-column coefficient estimate. At a coarse
level $h$, let $J_h$ map the coefficients of a trial in the chosen basis
to the coefficients of its selected level-$h$ polynomial. Write $C_h$
for the prescribed coefficient-column bound,
$Q_h$ for the ordinary norm bound of a normalized local trial, and
$E_h$ for its complete bound, including the complementary large-field
part. For the selected trajectory these bounds and all fixed coefficient
products are subpower in $a_h^{-1}$. The last-step endpoint estimate
and chessboard seminorm estimate for the residual required here are
\begin{align}
 \|J_h-I\|_{1\to1}&\le
           2(1+C_h)Q_hL_{h-1}^{-1/8}\le\tfrac14,
 \label{ch13:coarse-matrix}\\
 \|\text{trial residual}\|_{\mathcal C_k}
       &\le C_{k,m,d}E_h e^{-(\tau_h-\tau_m)/8}.
 \label{ch13:coarse-residual}
\end{align}
Here $k,m$ are fixed before $h$; $\mathcal C_k$ is the centered
cell tiling at level $k$ from Section~\ref{ch13:cells}. Both estimates must concern the history restricted from
the actual cutoff $N$, uniformly in $N\ge h$ and the physical period.

The first inequality is the endpoint comparison used in
Proposition~\ref{ch11:equation-comparison}. Its final $1/4$ is an
additional finite stage requirement: require $L_{h-1}^{1/8}\ge8(1+C_h)Q_h$. For the normalized small-field polynomial the local smooth cutoff
is identically one on the ordinary chart, so its ordinary derivative
constant is fixed before this block choice. Its global value is controlled by
the square-root large-field bound, by a larger coupling threshold as in
Section~\ref{ch11:equation-calibration}. This does not replace $E_h$ by
$Q_h$ in the second estimate.

The second inequality requires more than a small coefficient linear in the sources.
We prove its complete marked version next. The source weight is increased
across successive scales; it is not increased by its entire final factor
on an already formed coarse component.

\begin{lemma}[The coarse-trial residual on reflected tilings]
\label{ch13:coarse-residual-supplier}
On the selected histories, after the finite coupling-threshold requirements
specified below, \eqref{ch13:coarse-residual} holds with all incoming
regular and large-field terms involving at least two insertions retained. Its constant is independent
of the outer cutoff, the period, and the number of replica cells.
\end{lemma}
\begin{proof}
Fix the degree ceiling before the schedule. At level $i$, choose a
positive number $\zeta_i$ so that a term linear in the sources with complete normalized
size at most $\zeta_i$ fits every local attachment and adjacent
comparison bound of Chapter~\ref{ch11:insertions}. Include the finite
signed symmetry orbit in these bounds. More explicitly, if $D_i$ is
the complete adjacent comparison constant and $u_i>0$ the terminal
attachment bound, require
\[
 C_i\zeta_i\le1,\qquad 2D_i\zeta_{i+1}\le1,\qquad
 e^{81\sigma_i+162}Z_i^{\rm val}\zeta_i/u_i\le1.
\]
Here $Z_i^{\rm val}$ bounds the global section and its complementary
cutoff part; this symbol denotes a bound, not a partition function.
Also decrease $\zeta_i$ to put the incidence sum for terms linear in the sources below the
small constant in \eqref{ch11:higher-induction}. These choices concern
normalized coefficients, before the physical factor $a_i^{-d}$.

These positive upper bounds can be chosen with
$|\log\zeta_i|=O_j(1+\log\widehat\beta_i)$ on every fixed stage.
Indeed the section derivatives, inverse small-field thresholds, and
adjacent comparison constants have finite powers of
$\widehat\beta_i$, with stage-dependent constants. The terminal attachment bound may be taken as $u_i=\tfrac14\widehat\beta_i^{-1/2}$; the original large-field suppression
is reserved separately in the attachment estimate and need not be
inverted to define $u_i$. The factor $e^{81\sigma_i+162}$ is constant
on each stage. Taking the minimum of this finite list of reciprocal
bounds gives such a choice. At transitions between adjacent stages use the minimum
of both lists. All requirements below therefore follow by increasing
the coupling thresholds, since $B_i$ grows as a positive constant times
$(\log\widehat\beta_i)^2$ on a fixed stage.

Choose $0<\phi_i\le1$ such that
\begin{equation}
 \phi_i L_i^{d+1}\max(1,C_{i+1})\zeta_i\le\zeta_{i+1},
 \qquad \phi_i^{-1}\le e^{B_i/128}.
 \label{ch13:residual-height}
\end{equation}
For example take the minimum of one and the ratio in the first
inequality. The logarithm of its inverse is bounded by a finite
stage constant times $1+\log\widehat\beta_i$, plus
$(d+1)\log L_i$, and hence is eventually below $B_i/256$.
Any finite local incidence bound $\delta_i>0$ can simultaneously
be included so that $\delta_i\phi_i e^{B_i/128}\ge1$:
its inverse has at most another fixed power of
$\widehat\beta_i$. This verifies the comparison and attachment bounds as well as the displayed, weaker upper bound on
$\phi_i^{-1}$. These are countably many finite degree requests to
the single schedule of coupling thresholds.

Put $G(h,i)=e^{(\tau_h-\tau_i)/8}$ and choose
\[
 w=\frac{\zeta_m a_m^dG(h,m)}{2E_h},\qquad
 \omega_i^{\rm res}=w\prod_{j=m}^{i-1}\phi_j.
\]
The residual coefficients from \eqref{ch11:descent-error}, with the
fixed harmless column factors included in $E_h$, satisfy
$c_i+|x_i|\le E_h/G(h,i)$. Since the ratio of
$a_i^{-d}/G(h,i)$ at successive levels is $L_i^{d+1/8}$,
\eqref{ch13:residual-height} gives
\begin{equation}
 \omega_i^{\rm res} a_i^{-d}\bigl(c_i+|x_i|\bigr)
       \le \omega_i^{\rm res} a_i^{-d}E_h/G(h,i)\le\zeta_i/2
                  \quad(m\le i\le h).
 \label{ch13:residual-height-bound}
\end{equation}
In particular the negative normalized local insertion used in the subtraction has
normalized cost at most $\zeta_h$ at $h$. Above $h$, its inverse
coefficient bounds grow by at most $C_{i+1}$ per step; the same bound
recurrence controls these factors and the physical factor $L_i^d$.
Thus its terminal source satisfies the joint attachment bound at $N$,
and its complete push to $h$ is available at the weights $\omega_i^{\rm res}$.

Here is the support calculation for its terms involving at least two insertions, including those
formed above $h$. Open the centered cells at $m$ with spacing at least
$B_m^2$. This spacing persists because $L_i\ge3$ and the adjacent
support radii satisfy the schedule of \eqref{ch13:sharp-schedule}.
Apply reweighting on the incoming fine support $X$. If it contains
$r\ge2$ distinct marked ancestors, then
$B_{i+1}r\le64c_{B_{i+1}}(X)$ by \eqref{ch11:separated-cover}.
Since $B_i\le B_{i+1}$ and $\omega_i^{\rm res}/\omega_{i+1}^{\rm res}=\phi_i^{-1}$, one factor from the incoming exponential support weight gives
\[
 (\omega_i^{\rm res}/\omega_{i+1}^{\rm res})^r e^{-c_{B_{i+1}}(X)}
 \le e^{B_i r/128-c_{B_{i+1}}(X)}
 \le e^{-c_{B_{i+1}}(X)/2}.
\]
To make this argument iterable, let $A_i>4$ be the full exponent
maintained in the norm for terms involving at least two insertions at level $i$, and impose
\[
 10\cdot15^4(A_i+9)\le L_i(A_{i+1}-4).
\]
Reserve four units of the incoming $A_{i+1}$ norm for this weight
increase, site conversion, and attachments. The support contraction
of Chapter~\ref{ch11:insertions} converts the remaining exponent
$A_{i+1}-4$ into the full outgoing exponent $A_i$, including its local
reserves. Thus the estimate does not lose two exponent units at each
successive step. Choosing the bounds for terms linear in the sources small enough makes
\eqref{ch11:higher-induction} preserve the bound one in these full norms
for both arrays of terms involving at least two insertions. This calculation applies to every cardinality and
period. A large-field array is treated as one incoming factor with
its complete original prescribed set of large-field sites, not as a product of its coefficients linear in the sources. The terminal joint construction reserves the original
large-field decay for the complementary factors. Hence there is no
uncontrolled reweighting by $w^r$ at the intermediate level $h$.

At that level insert the positive local trial in every cell.
If $\mathscr A_h$ is the complete pushed negative source, the new polynomial is
\[
 \mathscr A_h\exp_{\rm nil}\Bigl(\sum_v z_vF_v\Bigr).
\]
In its local representation this adds $F_v$ only to the regular
entry linear in the sources. All regular terms involving at least two insertions and all original large-field
terms remain unchanged. The scalar part is exactly the coefficient of the identity in the global exponential.
The nonconstant part multiplied by the smooth cutoff has zero complementary profile because
its small-field restriction is included in its definition.
Linearity makes the new term linear in the sources precisely the residual whose
coefficients satisfy \eqref{ch13:residual-height-bound}; its adjacent
regular and large-field comparisons have bounds at most one by the
choice of $\zeta_i$. Descend this whole polynomial to $m$. Its regular terms linear in the sources have bound at most two, and the
large-field terms linear in the sources have bound at most one, and its two incidence bounds for terms involving at least two insertions are at most one.

Let $\mathsf P_{h,N}$ push densities from level $N$ to level $h$ by
successive applications of \eqref{ch04:block-operator} along the chosen
history. Its adjoint $\mathcal P_{N,h}$ pulls observables back from $h$ to $N$.
Let $\mathcal I$ index the cells of the reflected tiling and let $-U_v$
denote the terminal insertion used for cell $v$. In the product above
the coefficient of $\prod_{v\in\mathcal I}z_v$ at level $h$ is
\[
 \sum_{A\subseteq\mathcal I}
 \mathsf P_{h,N}\left(\rho_N\prod_{v\notin A}(-U_v)\right)
                       \prod_{v\in A}F_v.
\]
Duality and the disjoint-cone identity
\eqref{ch13:cone-factorization} identify its integral with the original
partition function times the expectation of
$\prod_v(\mathcal P_{N,h}F_v-U_v)$.
No product of separately pushed densities for single insertions has replaced it.
Equation~\eqref{ch11:root-bound} now yields
\[
 \|\mathcal P_{N,h}F-U\|_{\mathcal C_k}
 \le w^{-1}\exp\bigl((a_k/a_m)^4q_m^{\rm site}+5\bigr)
 =\frac{2e^{(a_k/a_m)^4q_m^{\rm site}+5}}{\zeta_ma_m^d}
       E_h e^{-(\tau_h-\tau_m)/8}.
\]
The prefactor depends only on the fixed stopping data. This proves
\eqref{ch13:coarse-residual}, once the cones have the room verified
below. All identities and estimates use the original selected history.
\end{proof}

\begin{proposition}[Coarse bounded approximation]
\label{ch13:bounded-singleton}
Use the selected histories with the stage requirements above. Fix a scaling-dimension bound $d$, a real reference coefficient $c$, a
sufficiently fine centered tiling $k$, and $\delta>0$. There are real
observables $\mathcal B_{N,p}(c,y)$ with a finite bound independent of
$N,p,y$, such that
\[
 \|\mathcal B_{N,p}(c,y)-O_{N,p}(d,c;y)\|_{\mathcal C_k(y)}\le\delta
\]
on the ordered cutoff and volume tails. Both observables are interior
to the same cell. The construction includes their signed and conjugated
images with the same bound.
\end{proposition}
\begin{proof}
Fix $m\ge k$ with $B_m^2\le a_k/a_m$. At a finer level $h$, use the
coarse trial
\[
 V_h(v;U)=a_h^{-d}\left(v_0+
       \mathbf1_{\text{small-field stencil}}
                    P_d^{\rm loc}[S_{1/t_h}v^\circ](U)\right).
\]
The scalar part is retained exactly. The full inspected stencil includes
the links needed by the relaxed small-field condition. For fixed $h$ this is
a bounded measurable function, with bound uniform in the outer cutoff
and period: its coefficient range and allowed couplings are compact,
and its section is a fixed polynomial in compact group variables.

Let $\mathcal K_{N,h}(U,dV)$ be the composition of the one-step probability kernels
$k^{(L_i)}_\beta(U,V)\,dV$ of \eqref{ch04:kernel}, in order from
$i=N-1$ down to $h$ and with the coupling of each step, and set
$\mathcal P_{N,h}F(U)=\int F(V)\mathcal K_{N,h}(U,dV)$.
It preserves real values and contracts the supremum norm. Locality here
must include every random variable used to choose a gauge-covariant block average or gauge
frame. Conditional on the preceding link field, use independent primitive
random variables for those choices and for each coarse-link sample.
The full ancestor cone of a stencil contains its fine-link reads and
all these random variables at every intermediate level. For disjoint
cones, conditioning and integrating the independent variables gives
\begin{equation}
 \prod_C\mathcal P_{N,h}F_C
                   =\mathcal P_{N,h}\left(\prod_CF_C\right).
 \label{ch13:cone-factorization}
\end{equation}
This does not assert that coarse links remain conditionally independent
after a shared frame choice has been averaged out. The full cones,
including that choice, must be disjoint. Here is an explicit uniform
radius bound. Let $j_i$ be the smoothing order in the step $i+1\to i$.
In Chapter~\ref{ch04:blocking}, each endpoint frame reads a box of
fine-lattice radius $(j_i+3)L_i$; the other endpoint is at distance
$L_i$, and the connector paths stay in these boxes and the intervening
strip. Thus every input link and every shared frame seed for an output
link lies within physical distance $(j_i+7)a_i$ of its initial endpoint,
in the maximum norm. This bound also includes eligibility conditions and
one boundary-link margin. If the coarse inspected stencil has radius
$r_d$ in level-$h$ units, its full ancestor cone therefore has radius at
most
\[
 r_da_h+\sum_{i=h}^{N-1}(j_i+7)a_i.
\]
The growing smoothing orders require a schedule condition, not merely
finiteness at each step. Impose $(j_i+7)\le a_i^{-1/2}$ at all sufficiently
late levels. This is a permitted finite coupling-threshold requirement: if
$\jmath_q$ is the largest smoothing order used at stage $q$ or its
permitted adjacent pairs, require
$b_q\log3\ge2\log(\jmath_q+7)$ among the coupling-threshold conditions of
Section~\ref{ch10:schedule}. When $J(i+1)=q>1$, that schedule has
$i\ge b_q$ and $a_i\le3^{-i}$, giving the desired inequality.
The finitely many initial levels are irrelevant after increasing $h$.
Since $a_{i+1}\le a_i/3$, the radius is bounded, independently of $N$, by
\[
 r_da_h+\frac{a_h^{1/2}}{1-3^{-1/2}}\longrightarrow0.
\]
Choose $h$ so that this is smaller than the fixed interior margin of
each reflected level-$k$ cell. The cones, including their shared seeds,
then lie strictly inside disjoint cells; sufficiently large physical
periods prevent wrapping identifications.

The bound \eqref{ch13:coarse-matrix} gives
$\|J_h^{-1}\|_1\le2$. Thus the reference insertion with coefficient
$c$ is a linear combination of the trial columns, with sum of absolute
coefficients bounded by
$2\|c\|_1\exp(\sum_{i<h}\log C_{i+1})$. Replace these columns by
$\mathcal P_{N,h}V_h(e_j)$ and call the resulting observable
$\mathcal B'$. Its supremum is finite for fixed $h$ and independent
of $N,p,y$.

Replicate the discrepancies before applying the kernel. At level $h$
the coarse factors multiply the complete source density, including its
terms involving at least two insertions. Equation~\eqref{ch13:cone-factorization} identifies the
resulting residual with the product of the pullback discrepancies.
Lemma~\ref{ch13:coarse-residual-supplier} and recombination give
\[
 \|\mathcal B'-O(c)\|_{\mathcal C_k}
 \le C_{k,m,d,c}\exp\left(\sum_{i<h}\log C_{i+1}\right)
                       E_h e^{-(\tau_h-\tau_m)/8}.
\]
The prefactor is $e^{o(\tau_h)}$, so choose $h$ to make this at most
$\delta$, before choosing the cutoff. For real $c$ take
$\mathcal B=\operatorname{Re}\mathcal B'$; conjugation invariance
and the triangle inequality for the chessboard seminorm preserve the
bound. Translations and the finite symmetry orbit use the same stencils
and constants. Outside the ordered tails set the approximants to zero
if a bound for every index is desired.
\end{proof}

\begin{corollary}[Real word differences]
\label{ch13:bounded-word}
On these selected histories, the bounded-real approximation used for phase
reflection in Section~\ref{ch15:phase-symmetry} holds on a cofinal set of
generic grids. Its supremum bound is independent of copy number and copy
position; approximation thresholds may depend on each fixed pair.
\end{corollary}
\begin{proof}
Choose the grid level after the degree bound and any prescribed lower
level, then fix a generic real word $w$ of length $n$. Include all its
signed and conjugated labels in a finite family. Choose centered cells
with margins fitting between its distinct coordinate projections, and
let $C$ bound the single-insertion seminorms there. Apply
Proposition~\ref{ch13:bounded-singleton} with error $\delta\le1$.
The approximants have seminorm at most $C+\delta$, regardless of their
possibly much larger supremum norms.

For a product of $q$ such separated factors, telescoping and the
chessboard inequality give
\[
 \left|\left\langle\prod_{a=1}^q\mathcal B_a-\prod_{a=1}^qO_a
               \right\rangle\right|
                  \le q\delta(C+\delta)^{q-1}.
\]
Form the difference $g_j$ of the bounded products for a translated word
and its linear time reflection, using the same single-insertion approximation
at every translated copy. If their supremum bounds for single insertions are $M_i$,
then $|g_j|\le2\prod_iM_i=:R_\varepsilon$, independent of all copy
positions and indices. For two sufficiently widely separated copies,
expand their product of differences into four products of $2n$ factors.
Choose $\delta$ with $8n\delta(C+1)^{2n-1}\le\varepsilon$.
The preceding inequality gives the desired pair error. Only finitely
many configurations are involved at each fixed pair, so ordered tails
can be intersected. The empty word difference is zero. Increasing the
initial lower-level request gives the required cofinal grids.
\end{proof}
